\documentclass[11pt]{article}
\usepackage{amsmath,amssymb,amsthm,bm,mathtools}
\usepackage{geometry}
\usepackage{hyperref}
\geometry{margin=1in}

\title{On the Link Between Bayesian Physics-Informed Neural Networks (B-PINNs) and Gaussian Processes (GPs)}
\author{}
\date{}

\newtheorem{proposition}{Proposition}
\newtheorem{corollary}{Corollary}
\newtheorem{remark}{Remark}

\begin{document}
\maketitle

\section{Introduction and Motivation}
Physics-Informed Neural Networks (PINNs) have emerged as a powerful paradigm for solving forward and inverse problems involving Partial Differential Equations (PDEs). By embedding physical laws directly into the loss function, PINNs can learn solutions to PDEs from scarce, scattered data. However, standard PINNs rely on deterministic optimization (e.g., Adam, L-BFGS) to converge on a single point estimate (Maximum a Posteriori, or MAP) of the network weights. These deterministic point-estimates are critically limited: they lack a principled statistical mechanism for \emph{uncertainty quantification} (UQ), leaving practitioners unable to trust the model's predictions in data-poor or highly noisy regimes.

To address this, Bayesian PINNs (B-PINNs) were introduced. A B-PINN places a prior distribution (typically isotropic Gaussian) over the weights and biases of the neural network, treating the PDE constraints as a soft likelihood penalty. Bayesian inference techniques, such as Hamiltonian Monte Carlo (HMC) or Variational Inference (VI), are then employed to sample from the posterior distribution of the network parameters. 

Despite their theoretical appeal, finite-width B-PINNs are notoriously challenging to deploy. The posterior landscape is highly non-convex and immensely high-dimensional. Sampling requires delicate tuning of HMC hyper-parameters (such as integration step sizes $\delta t$ and leapfrog trajectory lengths $L$) to prevent the Markov Chain from diverging or collapsing into localized energy wells. Furthermore, finite-width architectures can impose pathologies on the prior distributions of the predicted derivatives, threatening the stability of the physics gradients.

This motivates a fundamental theoretical investigation: \emph{What is the exact mathematical object that a Bayesian PINN approximates as its architecture scales?}

This study bridges deep Bayesian learning and classical spatial statistics. We exploit the established equivalence between infinite-width Bayesian Neural Networks (BNNs) and Neural Network Gaussian Processes (NNGPs). The core thesis of this work is:
\begin{center}
\textbf{In the infinite-width limit, a BNN prior induces an exact Gaussian Process (GP) prior on the surrogate solution field $u(\cdot)$.} \\
\vspace{0.2cm}
\textbf{If the PDE and boundary condition constraints enter linearly with respect to $u$, the entire B-PINN architecture reduces to a closed-form Gaussian Process regression problem formulated over linear operator observations.}
\end{center}
This equivalence allows us to derive an exact, analytical \emph{Physics-Informed Gaussian Process (PI-GP)} posterior. This analytical posterior completely bypasses the need for iterative neural network training and computationally unstable HMC sampling. More importantly, it serves as the absolute mathematical ground-truth. We can therefore systematically benchmark finite-width B-PINNs, quantifying their geometric convergence to the true posterior distribution as network width increases.

\section{The General B-PINN Formulation}

Let $u(x)$ be the unknown classical solution to a PDE defined on a spatial domain $D \subset \mathbb{R}^d$ with boundary $\Gamma$:
\begin{equation}
\mathcal{N}_x(u;\lambda)=f(x) \quad (x\in D),\qquad \mathcal{B}_x(u;\lambda)=b(x) \quad (x\in \Gamma),
\end{equation}
where $\mathcal{N}_x$ is a differential operator parameterized by the underlying physics scalars $\lambda$, and $\mathcal{B}_x$ is a boundary operator.

A B-PINN approximates $u(x)$ using an $L$-layer neural network surrogate $\tilde u(x;\theta)$, subject to parameters $\theta$. The implicit physics and boundary predictions are evaluated via automatic differentiation:
\begin{equation}
\tilde f(x;\theta)=\mathcal{N}_x(\tilde u(x;\theta);\lambda),\qquad \tilde b(x;\theta)=\mathcal{B}_x(\tilde u(x;\theta);\lambda)
\end{equation}

Assuming independent zero-mean Gaussian measurement noise, the joint likelihood function evaluating the fidelity of the parameters $\theta$ against the observed datasets $\mathcal{D}_u, \mathcal{D}_f, \mathcal{D}_b$ is:
\begin{equation}
P(\mathcal D\mid \theta)\propto \exp\left(
-\frac{1}{2\sigma_u^2}\sum_{x \in X_u} | \tilde u(x;\theta)-\bar u_x|^2
-\frac{1}{2\sigma_f^2}\sum_{x \in X_f}| \tilde f(x;\theta)-\bar f_x|^2
-\frac{1}{2\sigma_b^2}\sum_{x \in X_b}| \tilde b(x;\theta)-\bar b_x|^2
\right)
\end{equation}
where $\sigma_u, \sigma_f, \sigma_b$ define the respective measurement noise standard deviations. 

Coupled with a Gaussian prior on the network parameters $P(\theta) = \mathcal{N}(0, \sigma_\theta^2 I)$, the Bayesian framework defines the unnormalized posterior probability distribution of the parameters:
\begin{equation}
P(\theta\mid \mathcal D) = \frac{P(\mathcal D\mid \theta)\,P(\theta)}{P(\mathcal D)} \propto P(\mathcal D\mid \theta)\,P(\theta)
\end{equation}
The quantity $U(\theta) := -\log P(\theta\mid \mathcal D) \propto -\big[\log P(\mathcal D\mid \theta) + \log P(\theta)\big]$ forms the potential energy landscape traversed by the HMC sampler. Crucially, B-PINNs aim to evaluate the full posterior predictive integral over $u(x)$, rather than seeking the singular point-estimate $\operatorname*{argmin}_\theta U(\theta)$ analogous to standard PINN loss minimization.

\section{The Infinite-Width Limit: BNNs as NNGPs}

Consider a fully-connected multilayer perceptron (MLP) mapping inputs $x \in \mathbb{R}^{d_{in}}$ to a scalar output $u_L(x) \in \mathbb{R}$, featuring $L$ hidden layers of width $N_l$. The network computes $h^{(l)}(x) = \phi\big(W^{(l)} h^{(l-1)}(x) + b^{(l)}\big)$, where $W^{(l)}_{ij} \sim \mathcal{N}(0, \sigma_w^2 / N_{l-1})$ and $b^{(l)}_i \sim \mathcal{N}(0, \sigma_b^2)$.

\begin{proposition}[Convergence to the NNGP (Neal 1996, Lee et al. 2018)]
For any finite set of inputs $\{x_1,\dots,x_m\}$, as the widths of all hidden layers approach infinity ($N_l \to \infty$), the pre-activations of each layer converge in distribution to a multivariate Gaussian by the Central Limit Theorem. Consequently, the output surrogate function converges to a Gaussian Process:
\begin{equation}
u(\cdot;\theta) \xrightarrow[\min(N_1, \dots, N_L) \to \infty]{Distribution} \mathcal{GP}\big(0, k_{NNGP}(\cdot, \cdot)\big)
\end{equation}
The prior covariance function (kernel) $k_{NNGP}$ is computed deterministically via a layer-wise recursive formulation:
\begin{equation}
K^{(0)}(x, x') = \frac{\sigma_w^2}{d_{in}} x \cdot x' + \sigma_b^2
\end{equation}
\begin{equation}
K^{(l)}(x, x') = \sigma_w^2 \mathbb{E}_{(z, z') \sim \mathcal{N}(0, \Sigma^{(l-1)})} \Big[ \phi(z) \phi(z') \Big] + \sigma_b^2, \quad \text{where} \quad \Sigma^{(l-1)} = \begin{pmatrix} K^{(l-1)}(x, x) & K^{(l-1)}(x, x') \\ K^{(l-1)}(x', x) & K^{(l-1)}(x', x') \end{pmatrix}
\end{equation}
\end{proposition}

\section{Gaussian Processes under Linear Differential Operators}

Having analytically established that the B-PINN prior is $u \sim \mathcal{GP}\big(0, k_{NNGP}(x, x')\big)$, the conceptual bridge to extracting the physical constraints inherently concerns the behavior of differential operators across the probability measure.

\begin{proposition}[Linear Operator Closure of GPs]
Let $u(x)$ be a Gaussian Process distributed as $\mathcal{GP}\big(0, k(x, x')\big)$, and let $\mathcal{L}_x$ be a linear differential operator with respect to the spatial coordinates $x$. Assuming $k(x, x')$ possesses sufficient continuous mixed partial derivatives, the application of $\mathcal{L}_x$ to $u(x)$ yields another Gaussian Process. Furthermore, the joint process $[u, \mathcal{L}u]^T$ is a multivariate Gaussian Process tightly coupled by exact cross-covariances.
\end{proposition}

\begin{proof}
By the core definition of a Gaussian Process, any finite collection of function evaluations forms a multivariate Normal distribution. A linear differential operator $\mathcal{L}_x$ (e.g., the Laplacian $\Delta$, or directional derivative $\partial_x$) can be defined mathematically as the limit of a finite linear combination of function evaluations (i.e., finite difference approximations). Since the family of multivariate Gaussian distributions is closed under affine transformations and limits, the resulting field $\mathcal{L}_x u(x)$ is necessarily a highly-correlated Gaussian Process.

To explicitly construct the joint distribution, we compute the covariance matrix, inherently utilizing the linearity of expectation to commute the spatial differential operators outside the expectation bracket (justified by Dominated Convergence under sufficient kernel smoothness):
\begin{align}
\text{Cov}\big(\mathcal{L}_x u(x), \mathcal{L}_{x'} u(x')\big) &= \mathbb{E}\big[\mathcal{L}_x u(x) \cdot \mathcal{L}_{x'} u(x')\big] \nonumber \\
&= \mathcal{L}_x \mathcal{L}_{x'} \mathbb{E}\big[u(x) u(x')\big] \nonumber \\
&= \mathcal{L}_x \mathcal{L}_{x'} k(x, x') \label{eq:physics_covariance}
\end{align}
Similarly, the targeted cross-covariance connecting the spatial evaluation against the physical derivative evaluation is cleanly derived:
\begin{equation}
\text{Cov}\big(u(x), \mathcal{L}_{x'} u(x')\big) = \mathbb{E}\big[u(x) \cdot \mathcal{L}_{x'} u(x')\big] = \mathcal{L}_{x'} \mathbb{E}\big[u(x) u(x')\big] = \mathcal{L}_{x'} k(x, x') \label{eq:cross_covariance}
\end{equation}
\end{proof}

\section{The Exact Physics-Informed GP (PI-GP)}

If the underlying PDE physics constraint $\mathcal{N}_x(u)=f$ and boundary constraint $\mathcal{B}_x(u)=b$ are \emph{linear operators} governing $u$, we can formalize the B-PINN mathematically as an exact analytically solvable GP formulation, abandoning the artificial likelihood penalty scaling parameters.

Let $\mathcal{L}$ and $\mathcal{B}$ be linear. We stack our sets of noisy sensor measurements into an overarching block observation vector $Y \in \mathbb{R}^{N_u + N_f + N_b}$:
\begin{equation}
Y = \begin{bmatrix} \bar u \\ \bar f \\ \bar b \end{bmatrix} = \begin{bmatrix} u(X_u) \\ \mathcal{L}u(X_f) \\ \mathcal{B}u(X_b) \end{bmatrix} + \begin{bmatrix} \epsilon_u \\ \epsilon_f \\ \epsilon_b \end{bmatrix} = \mathcal{A} u + \epsilon, \quad \epsilon \sim \mathcal{N}(0, \Sigma_{\text{noise}})
\end{equation}
where the block noise covariance matrix aggregates over independent sensors: $\Sigma_{\text{noise}} = \text{blockdiag}(\sigma_u^2 I_{N_u}, \sigma_f^2 I_{N_f}, \sigma_b^2 I_{N_b})$.

By Proposition 2, the observations $Y$ follow a vast joint multivariate Normal distribution governed strictly by the structural derivatives of the $k(x, x')$ kernel matching equations \eqref{eq:physics_covariance} and \eqref{eq:cross_covariance}. We construct the dense generalized block covariance $\mathbf{K}$:
\begin{equation}
\mathbf{K} = \begin{bmatrix}
k(X_u, X_u) & \mathcal{L}_{x'}k(X_u, X_f) & \mathcal{B}_{x'}k(X_u, X_b) \\
\mathcal{L}_x k(X_f, X_u) & \mathcal{L}_x \mathcal{L}_{x'} k(X_f, X_f) & \mathcal{L}_x \mathcal{B}_{x'} k(X_f, X_b) \\
\mathcal{B}_x k(X_b, X_u) & \mathcal{B}_x \mathcal{L}_{x'} k(X_b, X_f) & \mathcal{B}_x \mathcal{B}_{x'} k(X_b, X_b)
\end{bmatrix}
\end{equation}

For any arbitrary spatial query point $x^*$, we construct the identical targeted cross-covariance row vector mapping the query to the established sensors:
\begin{equation}
\mathbf{k}_{*} = \begin{bmatrix} k(x^*, X_u) & \mathcal{L}_{x'}k(x^*, X_f) & \mathcal{B}_{x'}k(x^*, X_b) \end{bmatrix}
\end{equation}

The exact infinite-width, physics-informed posterior probability field evaluated at $x^*$ is analytically given by the closed-form standard GP regression predictive mapping:
\begin{equation}
\mu(x^*) = \mathbf{k}_{*} (\mathbf{K} + \Sigma_{\text{noise}})^{-1} Y
\end{equation}
\begin{equation}
\Sigma(x^*) = k(x^*, x^*) - \mathbf{k}_{*} (\mathbf{K} + \Sigma_{\text{noise}})^{-1} \mathbf{k}_{*}^T
\end{equation}

\subsection{Example Application: 1D Poisson Equation}
Consider the specific linear 1D Poisson Equation scaled by a mobility constant $\lambda$:
\begin{equation}
\lambda \partial_{xx}^2 u(x) = f(x) \implies \mathcal{L}_x u = \lambda \frac{d^2}{dx^2} u
\end{equation}
The off-diagonal cross-covariance linking an interior data sensor evaluation $u(x_u)$ to a physical constraint evaluation $\mathcal{L}_{x'} u(x_f)$ localized at a collocation sensor is derived explicitly by tracing the scaling coefficient:
\begin{equation}
\text{Cov}\big(u(x_u), \mathcal{L}_{x'} u(x_f)\big) = \lambda \left[ \frac{\partial^2}{\partial x'^2} k(x, x') \right]_{x=x_u, x'=x_f}
\end{equation}
Similarly, the highly autocorrelated physics-physics block covariance mapping the gradient trajectory relationship between two disparate collocation sensors $x_{f_i}$ and $x_{f_j}$ is defined by applying the localized operator onto both coordinate dimensions:
\begin{equation}
\text{Cov}\big(\mathcal{L}_x u(x_{f_i}), \mathcal{L}_{x'} u(x_{f_j})\big) = \lambda^2 \left[ \frac{\partial^4}{\partial x^2 \partial x'^2} k(x, x') \right]_{x=x_{f_i}, x'=x_{f_j}}
\end{equation}

By implementing these exact symbolic derivatives of the NNGP recursion function directly into GP regression, we generate the precise overarching probability density contours against which empirical B-PINNs converge as they scale.

\section{Summary \& Extension to Nonlinear Regimes}
\begin{itemize}
\item Bayesian Physics-Informed Neural Networks model the physical $u(\cdot)$ field using BNNs, projecting the PDE residuals back via an energetic likelihood function.
\item As layer width scales infinitely $N \rightarrow \infty$, the base surrogate prediction mathematically guarantees representation as a Neural Network Gaussian Process (NNGP) under Gaussian initializations.
\item If the acting partial differential operators are purely linear (e.g., $L_x = \lambda \Delta$), computing their covariance structurally reduces to computing mixed partial derivatives of the parent NNGP kernel.
\item Therefore, an infinitely wide B-PINN's Hamiltonian trajectory integral is rigorously identical—without any neural network sampling whatsoever—to the closed-form predictive bounds of a classically integrated block PI-GP.
\end{itemize}

What mathematical guarantees remain when addressing inherently non-linear regimes, such as the strongly advective $\mathcal{N}_x(u) = u \frac{\partial u}{\partial x}$ term characterizing the viscous Burgers' Equation?
While the governing Gaussian Process prior limit conceptually remains robust across deep architectures regardless of physics, applying a \emph{nonlinear operator} to a multivariate Gaussian shatters the probabilistic closure property. The joint structural distribution of $[u, \mathcal{N}_x(u)]^T$ is explicitly no longer Gaussian, preventing analytical mapping.

For nonlinear spaces, exploring exact closed-form limits is impossible. The theoretical state of the art instead relies on imposing generalized local approximations—such as constructing tangential Gaussian Processes linearized around the converged Maximum a Posteriori (MAP) field, or formalizing Laplace Approximations inside reproducing kernel Hilbert spaces. However, proving exact equivalencies across stable linear manifolds establishes the vital foundational milestone required before mapping finite architectures across turbulent boundaries.
\section{Empirical Evaluation: The Expressivity vs. Adaptivity Trade-off}

While the infinite-width limit provides a mathematically rigorous, closed-form posterior, empirical evaluation reveals a fundamental limitation of the Neural Network Gaussian Process (NNGP) in capturing highly oscillatory physical dynamics. We benchmarked the exact PI-GP formulation on the 1D Poisson equation where the true solution is $u(x) = \sin^3(6x)$. We tested two observational noise scales: $\sigma = 0.01$ and $\sigma = 0.1$.

\subsection{The Curse of the Stationary Kernel and "Lazy Learning"}
As shown in our empirical results (Figure \ref{fig:gp_comparison}), the exact PI-GP struggles to fit the target dynamics, and fails catastrophically when the measurement noise increases. To understand why the infinite-width limit fails where finite-width Deep Learning succeeds, we must examine the core assumption of the NNGP derivation.

\begin{figure}[h]
    \centering
    \includegraphics[width=\textwidth]{test_gp_results.png}
    \caption{Comparison of the exact PI-GP posterior mean and 2-sigma confidence interval against the true solution for the 1D Poisson equation with $u(x) = \sin^3(6x)$. (a) Low noise ($\sigma = 0.01$). (b) High noise ($\sigma = 0.1$).}
    \label{fig:gp_comparison}
\end{figure}

Taking the width $N \to \infty$ with i.i.d. Gaussian initializations invokes the Central Limit Theorem, which forces the network into the "lazy training" or Neural Tangent Kernel (NTK) regime. In this limit, the network's feature representations are strictly frozen. The resulting covariance function $k_{NNGP}(x, x')$ behaves as a rigid, stationary kernel that enforces a global length scale across the entire spatial domain. It cannot be highly oscillatory in one region and smooth in another.

\subsection{Noise-Induced Prior Dominance}
This rigidity mathematically explains the linear collapse observed at higher noise scales. In the predictive posterior mean equation:
\begin{equation}
\mu(x^*) = \mathbf{k}_{*} (\mathbf{K} + \Sigma_{\text{noise}})^{-1} Y
\end{equation}
the noise matrix $\Sigma_{\text{noise}}$ acts as Tikhonov regularization. When $\sigma = 0.1$, the noise variance scales up by a factor of $100$ compared to the $\sigma = 0.01$ case. Consequently, $\Sigma_{\text{noise}}$ heavily dominates the physics-physics block covariance $\mathcal{L}_x \mathcal{L}_{x'} k(X_f, X_f)$ inside the inversion. 

The GP effectively determines that the interior collocation points $X_f$ are too noisy to trust. Stripped of reliable interior data and unable to warp its feature space to isolate the high-frequency sine wave from the noise, the PI-GP falls back entirely on its zero-mean prior. It anchors only to the low-noise Dirichlet boundaries, resulting in an extreme low-pass filter that over-smooths the physics into a nearly linear interpolation.

\subsection{Representation Learning in Finite B-PINNs}
In stark contrast, literature demonstrates that finite-width B-PINNs (e.g., $N=50$) successfully regress this exact $\sin^3(6x)$ target even at the high $\sigma = 0.1$ noise scale, producing posterior means whose errors are strictly bounded by two standard deviations. 

Because finite networks do not fully trigger the Central Limit Theorem collapse, they retain the ability to perform \emph{representation learning}. During Hamiltonian Monte Carlo sampling, the parameters $\theta$ traverse the posterior landscape, actively adapting and warping the network's internal basis functions to align with the sharp local gradients of the underlying PDE. 

Therefore, while the infinite-width PI-GP provides perfect mathematical convexity and correct Gaussian priors for physical derivatives, finite-width B-PINNs remain practically superior for complex, highly oscillatory physics due to their capacity for non-stationary feature adaptation.
\end{document}